\section{引言}

% 写作提示：交代真实研究背景、统计研究对象、核心问题、数据能支持的推断范围及本文贡献；
% 避免“大数据时代”等泛化套话。下方背景和目标仅为结构示例，须按实际课题改写。

统计建模是数据分析的核心方法之一。通过建立统计模型，可以从数据中
提取有价值的信息和规律，为决策提供科学依据。

\subsection{研究背景}
在大数据时代，统计学方法在经济金融、生物医学、社会科学和工程技术等
各个领域都发挥着不可替代的作用。

\subsection{研究目标}
本文旨在通过统计建模方法，对给定的数据集进行全面分析，建立合理的统计模型。

% 技术路线图、研究框架图和算法流程图一律用 paper-diagram Skill 绘制——
% 优先套用其内置模板，产出可编辑 .drawio + PNG，
% 并用 check_layout.py / export_figure.py 校验后插入正文。
